\Askhsh{14820}{4}
\begin{ylh}{2}
    \item 2.1. Οι Πράξεις και οι Ιδιότητές τους
    \item 2.2. Διάταξη Πραγματικών Αριθμών
    \item 3.2. Η εξίσωση $x^{\nu}=a$
\end{ylh}

\begin{alist}
\item Να αποδείξετε ότι οι παρακάτω ανισότητες ισχύουν για κάθε $x \in \mathbb{R}$ και να βρείτε για ποιες τιμές του $x$ ισχύουν ως ισότητες.

\begin{rlist}
\item $x^2 + x + 1 \ge \dfrac{3}{4}$. \monades{4}
\item $x^2 - x + 1 \ge \dfrac{3}{4}$. \monades{4}
\end{rlist}
\item Να δείξετε ότι $(x^2 + x + 1)(x^2 - x + 1) > \dfrac{9}{16}$ για κάθε $x \in \mathbb{R}$. \monades{6}
\item Δίνεται η παράσταση $A = \dfrac{(x^3 - 1)(x^3 + 1)}{x^2 - 1}$.

\begin{rlist}
\item Να βρείτε για ποιες τιμές του $x \in \mathbb{R}$ ορίζεται η παράσταση $A$. \monades{5}
\item Με τη βοήθεια του β) ή με οποιοδήποτε άλλο τρόπο θέλετε, να εξετάσετε αν η παράσταση $A$ μπορεί να πάρει την τιμή $9/16$.
\monades{6}
\end{rlist}
\end{alist}

